\documentclass[12pt]{article}
\usepackage[utf8]{inputenc}
\usepackage[T1]{fontenc}
\usepackage[margin=1in]{geometry}
\usepackage{setspace}
\usepackage{xcolor}
\usepackage[colorlinks=true, urlcolor=blue!60!black, linkcolor=blue!60!black]{hyperref}
\setlength{\parindent}{2em}
\setlength{\parskip}{1.25em}
\renewcommand{\baselinestretch}{1.0}
\begin{document}
\begin{center}
{\LARGE \textit{Loper Bright Enterprises v. Raimondo}, 603 U.S. 369 (2024)\par}
\vspace{0.8em}
{\large Case Brief --- Four-Page Edition\par}
\end{center}
\vspace{1em}

\section{Memory Jogger}

Overruling \textit{Chevron}, the Court held 6--3 (6--2 in \textit{Loper Bright}) in a challenge to a rule making herring fishermen pay for federal observers that the APA requires courts to exercise independent judgment on whether an agency acted within its statutory authority, rather than deferring to an agency's reading of an ambiguous statute (\href{https://www.law.cornell.edu/supremecourt/text/22-451}{Cornell LII}).

\section{Detailed Case Facts}

The Magnuson-Stevens Fishery Conservation and Management Act, administered by the National Marine Fisheries Service (NMFS), lets fishery plans require vessels to carry federal ``observers.'' The statute expressly authorizes industry-funded observers for three groups---foreign vessels, limited access privilege programs, and the North Pacific (with fee caps)---but is silent on Atlantic herring. In 2020 NMFS approved an industry-funded monitoring program for the Atlantic herring fishery. When NMFS required an observer but did not supply one, owners had to hire certified observers at up to \$710 a day, cutting annual returns by up to 20\%.

\section{Procedural History}

Family-owned herring businesses sued under the MSA, which incorporates the APA, arguing that the statute did not authorize making fishermen pay. In \textit{Loper Bright}, the D.C. district court granted the government summary judgment, and a divided D.C. Circuit affirmed under \textit{Chevron} Step Two (Judge Walker dissenting). In \textit{Relentless}, the District of Rhode Island and the First Circuit also upheld the rule. The Supreme Court granted certiorari limited to whether \textit{Chevron} should be overruled, heard argument January 17, 2024, and decided June 28, 2024.

\section{Judicial Votes}

Vacated and remanded, \textbf{6--3} (\textit{Relentless}) and \textbf{6--2} (\textit{Loper Bright}; Justice Jackson recused).
\begin{itemize}
\item \textbf{Majority (6):} Chief Justice Roberts, joined by Justices Thomas, Alito, Gorsuch, Kavanaugh, and Barrett.
\item \textbf{Concurrences:} Justice Thomas; Justice Gorsuch.
\item \textbf{Dissent:} Justice Kagan, joined by Justice Sotomayor and, in \textit{Relentless}, Justice Jackson.
\end{itemize}

\section{Holding}

\textit{Chevron} is overruled. Under APA \S~706 and the Article III tradition of \textit{Marbury}, courts must exercise independent judgment on whether an agency acted within its statutory authority and may not defer simply because a statute is ambiguous. Courts may give \textit{Skidmore} ``respect'' to thorough, consistent agency judgments and must honor lawful delegations of discretion (to ``fill up the details'' or act as ``appropriate'' or ``reasonable''), policing their boundaries. Prior holdings that specific agency actions were lawful, including \textit{Chevron}'s own, remain protected by statutory stare decisis.

\section{Analysis of Opinions}

\textbf{Majority (Roberts).} \emph{Tradition and the APA.} Resolving legal ambiguity is ``the proper and peculiar province of the courts.'' From the Founding through the New Deal, courts respected executive constructions but kept independent judgment on law. Section 706 directs that ``the reviewing court shall decide all relevant questions of law,'' prescribing deference for policy and fact but none for law. \emph{\textit{Chevron}'s conflict.} \textit{Chevron}, which never cited the APA, made courts abandon their best reading for an agency's ``permissible'' one, even when the agency flip-flopped (\textit{Brand X}). Ambiguity is not implicit delegation: statutes have a ``single, best meaning,'' and agencies have no special competence in law. ``\textit{Chevron} does not prevent judges from making policy. It prevents them from judging.'' \emph{Stare decisis.} \textit{Chevron}'s reasoning was poor; ``ambiguity'' proved unworkable, producing a ``dizzying breakdance'' of exceptions that left a ``decaying husk''; and it undermined reliance by inviting reversals.

\textbf{Thomas (concurring).} \textit{Chevron} also violated the separation of powers by stripping courts of judicial power and handing judicial or legislative power to the executive.

\textbf{Gorsuch (concurring).} ``[T]he Court places a tombstone on \textit{Chevron} no one can miss.'' \textit{Chevron} tilted the scales toward government, undermined fair notice and lenity, and inflicted ``regulatory whiplash.''

\textbf{Kagan, with Sotomayor and Jackson (dissenting).} \textit{Chevron} sensibly reflected Congress's likely preference that expert, accountable agencies resolve technical gaps, such as whether a polymer is a ``protein'' or which squirrel populations are ``distinct.'' Section 706 is indeterminate, Congress legislated against \textit{Chevron} for forty years, and the majority replaced ``a rule of judicial humility'' with ``a rule of judicial hubris.''

\section{Examples: Future Applications}

\begin{enumerate}
\item \textbf{Same-Side: ``Emerging'' Compounds.} EPA, empowered to regulate ``emerging synthetic compounds,'' includes a chemical used since 1950. Whether it is ``emerging'' goes to the scope of authority, so the court decides \textit{de novo}.
\item \textbf{Same-Side: Barge Operators Paying for Inspectors.} A statute allows inspector fees in two programs, and the agency charges all operators. The court seeks the ``single, best meaning,'' which likely excludes fees elsewhere.
\item \textbf{Opposite-Side: Choice Within a Delegation.} EPA picks a 10-parts-per-billion threshold under a ``reasonable and necessary'' delegation. The court reviews only for reasoned decisionmaking.
\item \textbf{Opposite-Side: Rule Upheld Under \textit{Chevron}.} A company attacks a rule the Supreme Court upheld in 2010 under \textit{Chevron}. Statutory stare decisis protects the holding.
\item \textbf{Fence-Sitter: Agency Reversal.} A new administration repeals an FDA rule a circuit upheld under \textit{Chevron} Step Two and adopts the opposite reading. The old decision held only that the reading was permissible, and \textit{Loper Bright} did not say how it binds a contrary interpretation.
\end{enumerate}

\section{Critique}

\textbf{The ``single best meaning'' fiction.} Echoing Justice Kagan, critics say many technical phrases lack a fixed meaning, so \textit{Loper Bright} moves unavoidable policy choices from accountable agencies to life-tenured judges, risking circuit splits and ideological divergence.

\textbf{Delegation and stare decisis.} Because deference survives for delegations of discretion, disputes once fought at Step One have migrated to fights over which terms count as delegations. Courts also struggle to reconcile protection of prior \textit{Chevron} holdings with agencies' efforts to revise old rules now that \textit{Brand X} switching is gone.

\section{Key Quotations}

\begin{quote}
``\textit{Chevron} is overruled. Courts must exercise their independent judgment in deciding whether an agency has acted within its statutory authority, as the APA requires. Careful attention to the judgment of the Executive Branch may help inform that inquiry. And when a particular statute delegates authority to an agency consistent with constitutional limits, courts must respect the delegation, while ensuring that the agency acts within it. But courts need not and under the APA may not defer to an agency interpretation of the law simply because a statute is ambiguous.''
\end{quote}
\noindent --- Chief Justice Roberts, 603 U.S. at 412--413.

\begin{quote}
``In an agency case as in any other, though, even if some judges might (or might not) consider the statute ambiguous, there is a best reading all the same---`the reading the court would have reached' if no agency were involved. \textit{Chevron}, 467 U.S., at 843, n. 11. It therefore makes no sense to speak of a `permissible' interpretation that is not the one the court, after applying all relevant interpretive tools, concludes is best. In the business of statutory interpretation, if it is not the best, it is not permissible.''
\end{quote}
\noindent --- Chief Justice Roberts, 603 U.S. at 400.

\begin{quote}
``By doing so, however, we do not call into question prior cases that relied on the \textit{Chevron} framework. The holdings of those cases that specific agency actions are lawful---including the Clean Air Act holding of \textit{Chevron} itself---are still subject to statutory stare decisis despite our change in interpretive methodology.''
\end{quote}
\noindent --- Chief Justice Roberts, 603 U.S. at 412.

\end{document}
